% =========================================================
% 10. search_count() ORM Method
% =========================================================
\section{\texttt{search\_count()} ORM Method}

\begin{frame}[standout]
  \texttt{search\_count()} ORM Method
\end{frame}

\begin{frame}[fragile]{\texttt{SEARCH\_COUNT()} METHOD}
  \begin{itemize}
    \item The \texttt{search\_count()} method returns the \textbf{number of records}
          that match a domain.
    \item It is more efficient than searching all records and then using \texttt{len()}.
    \item The method returns an integer.
  \end{itemize}
  \begin{block}{Method Syntax}
\begin{lstlisting}
@api.model
def search_count(self, domain, limit=None):
    return super().search_count(domain, limit=limit)
\end{lstlisting}
    Important parts: \texttt{domain}, optional \texttt{limit}.
  \end{block}
\end{frame}

\begin{frame}[fragile]{BREAKING DOWN THE METHOD SIGNATURE}
  \begin{itemize}
    \item \texttt{@api.model}: Executed at the model level.
    \item \texttt{self}: The model, not a specific record.
    \item \texttt{domain}: Same domain language used by \texttt{search()}.
    \item \texttt{limit}: Optional maximum count (less commonly used).
  \end{itemize}
\begin{lstlisting}
domain = [('active', '=', True), ('age', '>=', 18)]
\end{lstlisting}
\end{frame}

\begin{frame}[fragile]{WHAT DOES \texttt{SEARCH\_COUNT()} RETURN?}
  \begin{itemize}
    \item It returns an \textbf{integer}.
  \end{itemize}
\begin{lstlisting}
total = self.env['student.student'].search_count([
    ('gender', '=', 'male'),
])
print(total)
\end{lstlisting}
  Output is (example):
\begin{lstlisting}
12
\end{lstlisting}
  \begin{itemize}
    \item No recordset is created.
    \item Only the count is computed.
  \end{itemize}
\end{frame}

\begin{frame}[fragile]{WHY NOT USE \texttt{LEN(SEARCH(...))}?}
  Less efficient pattern:
\begin{lstlisting}
total = len(self.env['student.student'].search([
    ('active', '=', True),
]))
\end{lstlisting}
  Better pattern:
\begin{lstlisting}
total = self.env['student.student'].search_count([
    ('active', '=', True),
])
\end{lstlisting}
  \begin{itemize}
    \item \texttt{search\_count()} avoids loading every matching record into Python.
    \item Important when tables are large.
  \end{itemize}
\end{frame}

\begin{frame}[fragile]{PRACTICAL EXAMPLES}
\begin{lstlisting}
adults = self.env['student.student'].search_count([
    ('age', '>=', 18),
])

archived = self.env['student.student'].search_count([
    ('active', '=', False),
])

has_email = self.env['student.student'].search_count([
    ('email', '!=', False),
])
\end{lstlisting}
  \begin{itemize}
    \item Useful for dashboards, KPIs, and validation checks.
  \end{itemize}
\end{frame}

\begin{frame}[fragile]{METHOD CALL}
  \textbf{Method Call}
\begin{lstlisting}
count = self.env['student.student'].search_count([
    ('name', 'ilike', 'ahmed'),
])
\end{lstlisting}
  \begin{itemize}
    \item Here, \texttt{.search\_count(...)} is the method call.
    \item Return type reminder: \textbf{int}.
  \end{itemize}
\end{frame}
